\documentclass[10pt,a4paper]{amsart}
\usepackage{amsmath,amssymb,amsthm,graphicx}
% Retained preamble for the Stacks Project chapter divisors.tex (Divisors).
% The chapter file pinned-original-source.tex carries only \input{preamble} on its first line, so no
% preamble text lives inside the chapter itself.  The Stacks Project uses one single shared preamble for
% every chapter; the snapshot stages it verbatim in sources/stacks-chow/preamble.tex.  The block below is
% copied mechanically, line for line, from that file's lines 14-18 (the Xy-pic setup) and lines 154-217
% (the native theorem declarations together with the mathematical macros).  Nothing was added, removed or
% redefined.  The wrapper class and the amsmath/amssymb/amsthm/graphicx base packages are supplied by the
% assembler, so the class/\externaldocument part of the shared preamble is not repeated here.
% For commutative diagrams we use Xy-pic
\usepackage[all]{xy}

% We use 2cell for 2-commutative diagrams.
\xyoption{2cell}
% Theorem environments.
%
\theoremstyle{plain}
\newtheorem{theorem}[subsection]{Theorem}
\newtheorem{proposition}[subsection]{Proposition}
\newtheorem{lemma}[subsection]{Lemma}

\theoremstyle{definition}
\newtheorem{definition}[subsection]{Definition}
\newtheorem{example}[subsection]{Example}
\newtheorem{exercise}[subsection]{Exercise}
\newtheorem{situation}[subsection]{Situation}

\theoremstyle{remark}
\newtheorem{remark}[subsection]{Remark}
\newtheorem{remarks}[subsection]{Remarks}

\numberwithin{equation}{subsection}

% Macros
%
\def\lim{\mathop{\mathrm{lim}}\nolimits}
\def\colim{\mathop{\mathrm{colim}}\nolimits}
\def\Spec{\mathop{\mathrm{Spec}}}
\def\Hom{\mathop{\mathrm{Hom}}\nolimits}
\def\Ext{\mathop{\mathrm{Ext}}\nolimits}
\def\SheafHom{\mathop{\mathcal{H}\!\mathit{om}}\nolimits}
\def\SheafExt{\mathop{\mathcal{E}\!\mathit{xt}}\nolimits}
\def\Sch{\mathit{Sch}}
\def\Mor{\mathop{\mathrm{Mor}}\nolimits}
\def\Ob{\mathop{\mathrm{Ob}}\nolimits}
\def\Sh{\mathop{\mathit{Sh}}\nolimits}
\def\NL{\mathop{N\!L}\nolimits}
\def\CH{\mathop{\mathrm{CH}}\nolimits}
\def\proetale{{pro\text{-}\acute{e}tale}}
\def\etale{{\acute{e}tale}}
\def\QCoh{\mathit{QCoh}}
\def\Ker{\mathop{\mathrm{Ker}}}
\def\Im{\mathop{\mathrm{Im}}}
\def\Coker{\mathop{\mathrm{Coker}}}
\def\Coim{\mathop{\mathrm{Coim}}}

% Boxtimes
%
\DeclareMathSymbol{\boxtimes}{\mathbin}{AMSa}{"02}

%
% Macros for moduli stacks/spaces
%
\def\QCohstack{\mathcal{QC}\!\mathit{oh}}
\def\Cohstack{\mathcal{C}\!\mathit{oh}}
\def\Spacesstack{\mathcal{S}\!\mathit{paces}}
\def\Quotfunctor{\mathrm{Quot}}
\def\Hilbfunctor{\mathrm{Hilb}}
\def\Curvesstack{\mathcal{C}\!\mathit{urves}}
\def\Polarizedstack{\mathcal{P}\!\mathit{olarized}}
\def\Complexesstack{\mathcal{C}\!\mathit{omplexes}}
% \Pic is the operator that assigns to X its picard group, usage \Pic(X)
% \Picardstack_{X/B} denotes the Picard stack of X over B
% \Picardfunctor_{X/B} denotes the Picard functor of X over B
\def\Pic{\mathop{\mathrm{Pic}}\nolimits}
\def\Picardstack{\mathcal{P}\!\mathit{ic}}
\def\Picardfunctor{\mathrm{Pic}}
\def\Deformationcategory{\mathcal{D}\!\mathit{ef}}

% --- appended 2026-10-12 (Chief): the project own definitions of the reference/slogan/history
% display environments, copied mechanically from the shared Stacks preamble lines 8-12
% (\usepackage{verbatim} plus the three \newenvironment lines, which the chapter preamble keeps
% verbatim from the project preamble). Nothing added, removed or redefined: these environments
% carry no mathematical content and the project itself typesets them as comments. Without these
% three lines a counted statement whose own text uses them fails with Environment slogan undefined.
\usepackage{verbatim}
\newenvironment{reference}{\comment}{\endcomment}
\newenvironment{slogan}{\comment}{\endcomment}
\newenvironment{history}{\comment}{\endcomment}
\usepackage{cleveref}
\crefname{originallemma}{original lemma}{original lemmas}
\Crefname{originallemma}{Original lemma}{Original lemmas}
\makeatletter
\newlabel{algebra-lemma-associated-primes-localize}{{algebra-lemma-associated-primes-localize}{source}{}{sourcepointer-1}{}}
\newlabel{algebra-lemma-associated-primes-localize@cref}{{[originallemma][0][]algebra-lemma-associated-primes-localize}{[1][0][]source}{}{}{}}
\makeatother
\crefname{originallemma}{original lemma}{original lemmas}
\Crefname{originallemma}{Original lemma}{Original lemmas}
\makeatletter
\newlabel{properties-lemma-locally-Noetherian}{{properties-lemma-locally-Noetherian}{source}{}{sourcepointer-2}{}}
\newlabel{properties-lemma-locally-Noetherian@cref}{{[originallemma][0][]properties-lemma-locally-Noetherian}{[1][0][]source}{}{}{}}
\makeatother
\crefname{originallemma}{original lemma}{original lemmas}
\Crefname{originallemma}{Original lemma}{Original lemmas}
\makeatletter
\newlabel{schemes-lemma-spec-sheaves}{{schemes-lemma-spec-sheaves}{source}{}{sourcepointer-3}{}}
\newlabel{schemes-lemma-spec-sheaves@cref}{{[originallemma][0][]schemes-lemma-spec-sheaves}{[1][0][]source}{}{}{}}
\makeatother
\title{Associated points over an affine open are detected on the module of sections by the corresponding prime}
\author{The Stacks Project Authors}
\date{stacks-project:divisors.tex:local-snapshot-20261012; GNU FDL 1.2 or later}
\begin{document}
\maketitle
\noindent\textit{Editorial note.} All delivered mathematical text is the author's own, copied from the pinned Stacks Project chapter file pinned-original-source.tex; no statement or proof step was rewritten, reordered or removed.\par\medskip
\noindent\textit{Editorial note.} Following the multi-claim convention for this source, source\_url is the source-level URL https://stacks.math.columbia.edu/ and the source label lemma-associated-affine-open is carried in claim\_key and here rather than in a per-claim tag URL.\par\medskip
\noindent\textit{Editorial note.} The preamble is the chapter's shared Stacks Project preamble, retained as sources/stacks-divisors/retained-preamble.tex: a line-for-line excerpt of the shared Stacks preamble staged with this snapshot. No macro was invented or redefined, and this card adds no preamble text of its own.\par\medskip
\noindent\textit{Editorial note.} Statement span 67-84 and proof span 86-117 come from the halfA pool card operations/pool-cards/stacks-divisors-halfA.json and were re-read in the pinned chapter file before this card was written.\par\medskip
\noindent\textit{Editorial note.} The proof opens with the author's sentence that the statement also follows from a cited algebra lemma, and then says that a direct argument can be given as well; both the citation and the direct argument are preserved, and the direct argument is what is counted.\par\medskip
\noindent\textit{Editorial note.} 3 cross-references occur inside the counted spans; they are listed in external\_references as navigation pointers only.\par\medskip
\noindent\textit{Allowed standard prerequisites (editorial).}\begin{itemize}
\item An associated prime of a module over a ring is the annihilator of some element, and localization is exact. Used at proof lines 88-94 and 97-107. Produces the element whose annihilator is the prime in the first direction, and the element of the localization whose annihilator is the localized prime in the second.
\item The stalk of a quasi-coherent module on an affine open at a prime is the localization of the module of sections at that prime, cited from the schemes chapter. Cited at proof lines 92-93. Identifies the localized module with the stalk so that association at the point follows from association at the prime.
\item For a finitely generated prime ideal contained in the annihilator up to localization there is an element outside the prime killing the quotient. Used at proof lines 104-108. Upgrades association of the localized element to association of an honest element of the module of sections.
\item In a locally Noetherian scheme the coordinate rings of affine opens are Noetherian, cited from the properties chapter. Cited at proof lines 112-114. Makes every prime finitely generated, which is the hypothesis of the converse direction.
\end{itemize}\medskip
\ifdefined\hypertarget\hypertarget{sourcepointer-1}{}\fi
\noindent\textit{Source reference.} Original lemma algebra-lemma-associated-primes-localize: sources/stacks-divisors/pinned-original-source.tex, source lines 88--88. Same-chapter reference used in the retained author text at source line(s) 88, 6303. It is supplied only as a navigation pointer so that the retained text compiles; the counted item delivers the author's own statement and the author's complete proof, and no part of the referenced item is copied in or substituted for the counted argument.\par\medskip
\ifdefined\hypertarget\hypertarget{sourcepointer-2}{}\fi
\noindent\textit{Source reference.} Original lemma properties-lemma-locally-Noetherian: sources/stacks-divisors/pinned-original-source.tex, source lines 115--115. Same-chapter reference used in the retained author text at source line(s) 115, 164, 850. It is supplied only as a navigation pointer so that the retained text compiles; the counted item delivers the author's own statement and the author's complete proof, and no part of the referenced item is copied in or substituted for the counted argument.\par\medskip
\ifdefined\hypertarget\hypertarget{sourcepointer-3}{}\fi
\noindent\textit{Source reference.} Original lemma schemes-lemma-spec-sheaves: sources/stacks-divisors/pinned-original-source.tex, source lines 95--95. Same-chapter reference used in the retained author text at source line(s) 95. It is supplied only as a navigation pointer so that the retained text compiles; the counted item delivers the author's own statement and the author's complete proof, and no part of the referenced item is copied in or substituted for the counted argument.\par\medskip
\begin{lemma}
\label{lemma-associated-affine-open}
Let $X$ be a scheme. Let $\mathcal{F}$ be a quasi-coherent sheaf on $X$.
Let $\Spec(A) = U \subset X$ be an affine open, and set
$M = \Gamma(U, \mathcal{F})$.
Let $x \in U$, and let $\mathfrak p \subset A$ be the corresponding prime.
\begin{enumerate}
\item If $\mathfrak p$ is associated to $M$, then $x$ is associated
to $\mathcal{F}$.
\item If $\mathfrak p$ is finitely generated, then the converse holds
as well.
\end{enumerate}
In particular, if $X$ is locally Noetherian, then the equivalence
$$
\mathfrak p \in \text{Ass}(M) \Leftrightarrow x \in \text{Ass}(\mathcal{F})
$$
holds for all pairs $(\mathfrak p, x)$ as above.
\end{lemma}

\begin{proof}
This follows from
Algebra, Lemma \ref{algebra-lemma-associated-primes-localize}.
But we can also argue directly as follows.
Suppose $\mathfrak p$ is associated to $M$.
Then there exists an $m \in M$ whose annihilator is $\mathfrak p$.
Since localization is exact we see that
$\mathfrak pA_{\mathfrak p}$ is the annihilator of
$m/1 \in M_{\mathfrak p}$. Since $M_{\mathfrak p} = \mathcal{F}_x$
(Schemes, Lemma \ref{schemes-lemma-spec-sheaves})
we conclude that $x$ is associated to $\mathcal{F}$.

\medskip\noindent
Conversely, assume that $x$ is associated to $\mathcal{F}$,
and $\mathfrak p$ is finitely generated.
As $x$ is associated to $\mathcal{F}$
there exists an element $m' \in M_{\mathfrak p}$ whose
annihilator is $\mathfrak pA_{\mathfrak p}$. Write
$m' = m/f$ for some $f \in A$, $f \not \in \mathfrak p$.
The annihilator $I$ of $m$ is an ideal of $A$ such that
$IA_{\mathfrak p} = \mathfrak pA_{\mathfrak p}$. Hence
$I \subset \mathfrak p$, and $(\mathfrak p/I)_{\mathfrak p} = 0$.
Since $\mathfrak p$ is finitely generated,
there exists a $g \in A$, $g \not \in \mathfrak p$ such that
$g(\mathfrak p/I) = 0$. Hence the annihilator of $gm$ is
$\mathfrak p$ and we win.

\medskip\noindent
If $X$ is locally Noetherian, then $A$ is Noetherian
(Properties, Lemma \ref{properties-lemma-locally-Noetherian})
and $\mathfrak p$ is always finitely generated.
\end{proof}

\end{document}
